%! Exponential Function Context and Data Modeling — Unit 4, Lesson 4.6 — Group Activity (Answer Key)
\documentclass[10pt]{article}
\usepackage{precalculus-article}
\usepackage{precalculus-key}   % pulls in -boxes; do NOT also load -boxes
\newcommand{\TallMath}[1]{$\displaystyle #1\rule[-1.4em]{0pt}{3.2em}$}

\begin{document}

\pageheader{Unit 4, Lesson 4.6}{Group Activity --- Answer Key}
\namepartnerperiod

\vspace{0.1in}

\begin{scenariobox}[Building Exponential Models]{sky}
\small
Exponential models appear in medicine, ecology, finance, and physics. In each tier below, you will
build and apply an exponential model, interpret its base, and use it to answer a contextual question.
Work with your group on the tier assigned by your teacher.
\end{scenariobox}

\vspace{0.1in}

\begin{tcolorbox}[colback=white, colframe=black!40,
    title=\textbf{Tier R --- Remediate},
    fonttitle=\bfseries, arc=2mm, left=3mm, right=3mm, top=2mm, bottom=2mm]

\textbf{Context:} A drug's concentration in the bloodstream halves every 4 hours. The initial
concentration is 200 mg/L.

\begin{enumerate}[itemsep=10pt]
  \item The model is $C(t) = 200 \cdot (0.5)^{t/4}$, where $t$ is hours. Complete the table.

        \vspace{0.06in}
        \renewcommand{\arraystretch}{1.8}
        \begin{tabular}{c|c}
        \textbf{Time $t$ (hours)} & \textbf{Concentration $C(t)$ (mg/L)} \\ \hline
        0  & \ans{200} \\
        4  & \ans{100} \\
        8  & \ans{50} \\
        12 & \ans{25} \\
        \end{tabular}

  \item What is the base of this exponential model? What does it tell you about the drug?

        \vspace{0.06in}
        Base: \ans{0.5} \quad Meaning: \ans{Every 4 hours, the concentration is multiplied by 0.5 (halved).}

  \item Use the model to find the concentration after 6 hours. Show your work.

        \vspace{0.06in}
        $C(6) = 200 \cdot (0.5)^{6/4} = $ \ans{$200 \cdot (0.5)^{1.5} = 200 \cdot 0.3536$} $\approx$ \ans{70.7} mg/L
\end{enumerate}
\end{tcolorbox}

\vspace{0.1in}

\begin{tcolorbox}[colback=white, colframe=black!40,
    title=\textbf{Tier A --- Apply},
    fonttitle=\bfseries, arc=2mm, left=3mm, right=3mm, top=2mm, bottom=2mm]

\textbf{Context:} A town's population was 4{,}000 in Year 2 and 6{,}250 in Year 5.
Assume exponential growth.

\begin{enumerate}[itemsep=10pt]
  \item Set up and solve a system of equations to find $a$ and $b$ in $P(t) = ab^t$.

        \vspace{0.06in}
        Equation 1: \ans{$4000 = ab^2$}

        Equation 2: \ans{$6250 = ab^5$}

        Divide to find $b$: \ans{$\dfrac{6250}{4000} = b^{5-2} = b^3$, so $b^3 = 1.5625$}

        \vspace{0.04in}
        $b = $ \ans{$1.5625^{1/3} \approx 1.16$} \quad then $a = $ \ans{$4000 / (1.16)^2 \approx 2973$}

        \vspace{0.04in}
        Model: $P(t) = $ \ans{$2973 \cdot (1.16)^t$}

  \item What is the annual growth rate as a percent? \ans{Approximately 16\% per year ($b \approx 1.16$, so $r \approx 0.16$).}

  \item Predict the population at Year 10. Then note: is this prediction within a reasonable domain?

        \vspace{0.06in}
        $P(10) = $ \ans{$2973 \cdot (1.16)^{10}$} $\approx$ \ans{13{,}100}

        \vspace{0.04in}
        Domain note: \ans{Year 10 is close to the data range; extrapolating far beyond Year 5 requires caution.}
\end{enumerate}
\end{tcolorbox}

\vspace{0.1in}

\begin{tcolorbox}[colback=white, colframe=black!40,
    title=\textbf{Tier E --- Extend},
    fonttitle=\bfseries, arc=2mm, left=3mm, right=3mm, top=2mm, bottom=2mm]

\textbf{Context:} $f(d) = 3 \cdot 2^d$ models the daily growth of a bacterial colony, where
$d$ is the number of days and $f(d)$ is the colony size in thousands.

\begin{enumerate}[itemsep=10pt]
  \item Rewrite $f(d)$ as a model $g(w) = a \cdot b^w$, where $w$ is the number of \textbf{weeks}.
        Show the algebra clearly.

        \vspace{0.06in}
        $g(w) = $ \ans{$3 \cdot 2^{7w} = 3 \cdot (2^7)^w = 3 \cdot 128^w$} \quad What does the base $b$ mean in this context?

        \vspace{0.04in}
        \ans{The colony multiplies by a factor of 128 every week.}

  \item Rewrite $f(d)$ as a model $h(n) = a \cdot c^n$, where $n$ is the number of \textbf{3-day
        periods} (so $d = 3n$). What does the base $c$ mean?

        \vspace{0.06in}
        $h(n) = $ \ans{$3 \cdot 2^{3n} = 3 \cdot (2^3)^n = 3 \cdot 8^n$} \quad Meaning of base $c$: \ans{The colony multiplies by 8 every 3 days.}

  \item After how many days does the colony first exceed 10{,}000 thousand (i.e., $f(d) > 10000$)?

        \vspace{0.04in}
        \renewcommand{\arraystretch}{1.6}
        \begin{tabular}{c|c}
        $d$ & $f(d) = 3 \cdot 2^d$ \\ \hline
        10 & \ans{3072} \\
        11 & \ans{6144} \\
        12 & \ans{12288} \\
        \end{tabular}

        \vspace{0.06in}
        The colony first exceeds 10{,}000 thousand after $d = $ \ans{12} days.
\end{enumerate}
\end{tcolorbox}

\begin{teachernote}
Tier R: $C(6) = 200 \cdot (0.5)^{1.5} = 200/2\sqrt{2} \approx 70.7$ mg/L. Accept 70--71 mg/L.

Tier A: Exact value of $b = (6250/4000)^{1/3} = (25/16)^{1/3}$. Numerically $\approx 1.1604$. Then $a = 4000/b^2 \approx 2973$. Students who get $b \approx 1.16$ and a reasonable $a$ should be credited.

Tier E: Item 3 --- $3 \cdot 2^{12} = 3 \cdot 4096 = 12288 > 10000$, and $3 \cdot 2^{11} = 6144 < 10000$, so the colony first exceeds 10{,}000 thousand after day 12.
\end{teachernote}

\end{document}
